\documentclass[level=1]{kidscode}
\newcommand{\lvA}{\faSeedling\ Just starting}
\newcommand{\lvB}{\faShoePrints\ Getting there}
\newcommand{\lvC}{\faCheckCircle\ Got it}
\newcommand{\lvD}{\faRocket\ Expert}
\begin{document}
\kctitle{My Interactive Art Show}{Final project brief, plan sheet and rubric (weeks 10--12)}

\section{Your mission}
You are an artist who paints with code. In weeks 10, 11 and 12 you will make a
Python turtle program that is a little \textbf{interactive art show}:

\begin{enumerate}[itemsep=2pt]
  \item it says hello to the person watching (the \enquote{viewer}),
  \item it asks them one or two questions,
  \item it draws a picture or pattern that changes because of their answers,
  \item and it has a surprise, so it is a bit different every time you run it.
\end{enumerate}

In week 12 you will show it to your family at the \textbf{Showcase} and explain
one part of your code. You can work alone or with your partner. If you work in a pair,
you both need to be able to explain every part.

\begin{tip}[An example]
Your teacher has a finished example called \emph{Night Sky Postcard}. It asks
\enquote{How many stars?} and \enquote{Day or night?}, then draws the sky,
a moon or a sun, random stars and a row of houses. Ask to see it, then make
something that is \textbf{yours}.
\end{tip}

\section{The six ingredients}
Your finished program must use all six. Tick each one when it is in your code and working.

\renewcommand{\arraystretch}{1.35}
\begin{tabularx}{\linewidth}{@{}p{0.6cm}>{\bfseries}p{2.6cm}X@{}}
\toprule
 & Ingredient & What it looks like in Python\\
\midrule
\tickbox & A question & \texttt{answer = input("How many fish? ")}\\
\tickbox & A variable & \texttt{fish\_count = int(answer)} \quad (a labelled box that remembers)\\
\tickbox & An \texttt{if} & \texttt{if mood == "happy":} \dots{} \texttt{else:} \dots\\
\tickbox & A loop & \texttt{for i in range(fish\_count):} \dots\\
\tickbox & Something random & \texttt{x = random.randint(-350, 350)}\\
\tickbox & My own function & \texttt{def fish(x, y, colour):} \dots{} and then calling it: \texttt{fish(0, 0, ORANGE)}\\
\bottomrule
\end{tabularx}

\section{Idea menu}
Pick one, mix two, or invent your own. Every idea can be an artwork or a little game.

\begin{multicols}{2}
\begin{description}[leftmargin=0pt,itemsep=4pt,font=\bfseries]
  \item[Underwater world] Asks \enquote{How many fish?} and \enquote{Deep or shallow?}. Your \texttt{fish(x, y, colour)} function swims in random places.
  \item[Flower garden] Asks \enquote{How many flowers?} and \enquote{Sunny or rainy?}. Rain draws lines, sun draws a sun.
  \item[Monster maker] Asks \enquote{How many eyes?} and \enquote{Happy or grumpy?}. The \texttt{if} picks the mouth shape.
  \item[Space trip] Asks \enquote{How many planets?}. Planets get random sizes; a \texttt{rocket()} function flies past.
  \item[Kaleidoscope] Asks \enquote{How many sides?}. A loop turns and repeats your shape with random colours.
  \item[Snow day] Asks \enquote{How much snow: a little or a lot?}. Snowflakes are made by a \texttt{snowflake()} function.
  \item[Fireworks show] Asks \enquote{How many fireworks?}. Each firework is a burst of lines in a random place.
  \item[Guess-the-shape game] The program secretly picks a random number of sides, draws the shape, asks \enquote{How many sides?} and the \texttt{if} says \enquote{Correct!} or \enquote{Try again}.
\end{description}
\end{multicols}

\begin{safety}[Your art, your privacy]
Use a nickname on your artwork, never your real full name, address or school.
Do not ask the viewer for private things either: ask about colours, numbers, moods or animals instead.
At the showcase we show projects on the screen only. Nothing is uploaded.
\end{safety}

\section{Where things go on the screen}
The turtle window is 800 steps wide and 600 steps tall. The middle is $(0, 0)$.
Use \texttt{t.penup()} and \texttt{t.goto(x, y)} to jump to a spot without drawing.

\begin{center}
\begin{tikzpicture}[x=0.012cm,y=0.012cm,font=\small]
  \fill[kcsky!25] (-400,-300) rectangle (400,300);
  \draw[kcgrey!50,step=100] (-400,-300) grid (400,300);
  \draw[thick] (-400,-300) rectangle (400,300);
  \draw[-{Latex},thick] (-400,0) -- (430,0) node[right] {$x$};
  \draw[-{Latex},thick] (0,-300) -- (0,330) node[above] {$y$};
  \fill (0,0) circle (6);
  \node[below right] at (0,0) {$(0, 0)$ middle};
  \node[below left] at (400,0) {$(400, 0)$};
  \node[below right] at (-400,0) {$(-400, 0)$};
  \node[below right] at (0,300) {$(0, 300)$ top};
  \node[above right] at (0,-300) {$(0, -300)$ bottom};
  \fill[kcorange] (200,150) circle (8);
  \node[right] at (210,150) {$(200, 150)$: right and up};
  \fill[kcpurple] (-250,-150) rectangle ++(16,16);
  \node[right] at (-230,-142) {$(-250, -150)$: left and down};
\end{tikzpicture}
\end{center}

\clearpage
\section{My plan sheet (Milestone 1)}
\namedateline
\textbf{My project's name:} \rule{9cm}{0.4pt}

\textbf{Sketch of my picture} (draw it here; label the parts):
\answerspace{5.5cm}

\textbf{What my program asks} (1 or 2 questions):
\answerlines{2}

\textbf{What it draws} (and what changes because of the answers):
\answerlines{2}

\textbf{My function}\quad Name: \rule{4cm}{0.4pt}\quad It needs: \rule{5cm}{0.4pt}\\[2pt]
{\small (for example \texttt{fish(x, y, colour)} needs a place and a colour)}\par
What it draws:
\answerlines{1}

\textbf{Where are my ingredients?} Write a few words for each.
\begin{checklist}
  \item My loop repeats: \rule{9cm}{0.4pt}
  \item My \texttt{if} decides: \rule{9cm}{0.4pt}
  \item My random surprise is: \rule{8cm}{0.4pt}
  \item My variable remembers: \rule{8cm}{0.4pt}
\end{checklist}
Teacher OK for my plan: \rule{4cm}{0.4pt}

\clearpage
\section{Milestones}
Build in small steps. Each milestone has a check. Tick it when it is done.

\renewcommand{\arraystretch}{1.4}
\begin{tabularx}{\linewidth}{@{}p{0.6cm}>{\bfseries}p{3.2cm}X@{}}
\toprule
 & Milestone & Done when\dots\\
\midrule
\tickbox & M1 Plan\newline{\normalfont\small end of week 10} & My plan sheet is filled in and my teacher said OK. I opened \texttt{my\_art\_show.py} and saved it with my project's name.\\
\tickbox & M2 Core drawing works\newline{\normalfont\small week 11, first half} & My function draws my main thing, and I call it at least twice in different places. The program runs without red errors.\\
\tickbox & M3 Surprise and polish\newline{\normalfont\small week 11, second half} & My questions, my \texttt{if}, my loop and my random surprise all work. A partner tested it and gave me \enquote{two stars and a wish}. I added a title or message.\\
\tickbox & M4 Showcase\newline{\normalfont\small week 12} & I can run my program for a visitor and explain one part of my code in my own words.\\
\bottomrule
\end{tabularx}

\begin{tip}[Small steps beat big leaps]
Change one thing, press \textbf{Run}, look. Then change the next thing.
If something breaks, you know it was the last thing you changed.
Save often with \keys{Cmd} + \keys{S}.
\end{tip}

\section{How my project is assessed}
Your teacher uses this rubric. Use it yourself too: circle where you are now, and again
at the end. Every column is a good place to be on the way.

{\small\renewcommand{\arraystretch}{1.3}
\begin{tabularx}{\linewidth}{@{}>{\bfseries\RaggedRight}p{2.5cm}*{4}{>{\RaggedRight\arraybackslash}X}@{}}
\toprule
I can\dots & \lvA & \lvB & \lvC & \lvD\\
\midrule
make it my own (creativity) & I copied an example. & I changed colours or sizes in an example. & My picture has my own idea, my own function and a surprise. & My idea is unique and I added extra details or a twist nobody expected.\\
\midrule
make my code work & It does not run yet. & It runs, but some parts are missing or broken. & It runs from start to end without errors. & It also copes with odd answers (like \enquote{lots} instead of a number).\\
\midrule
use the six ingredients & I used 1 or 2. & I used 3 or 4. & I used all 6. & I used all 6, and my function has parameters I really use (like \texttt{x}, \texttt{y}, \texttt{colour}).\\
\midrule
explain my code & I can say what my program does. & I can point to my loop and my function. & I can explain how one part of my code works, line by line. & I can explain why I chose it, and how I fixed a bug.\\
\bottomrule
\end{tabularx}}

\section{Need a different path?}
\begin{tip}[Support path: make it in Scratch]
You can build the same art show in Scratch instead. The rubric is exactly the same,
and you use the same six ingredients:
\begin{center}
\small\renewcommand{\arraystretch}{1.3}
\begin{tabular}{@{}ll@{}}
\toprule
Ingredient & Scratch block\\
\midrule
A question & \textbf{ask} \dots{} \textbf{and wait}, then the \textbf{answer} block\\
A variable & \textbf{Make a Variable}, then \textbf{set} / \textbf{change}\\
An if & \textbf{if} \dots{} \textbf{then} / \textbf{if} \dots{} \textbf{then} \dots{} \textbf{else}\\
A loop & \textbf{repeat} \dots{} or \textbf{forever}\\
Something random & \textbf{pick random} \dots{} \textbf{to} \dots\\
My own function & \textbf{Make a Block} (a My Block, with inputs if you can)\\
\bottomrule
\end{tabular}
\end{center}
\end{tip}

\begin{stretchchallenge}[Extension path: react to key presses]
Make your art show react while it is open. Put this \textbf{before} \texttt{turtle.done()}.
Pressing the space bar adds a new star in a random place:
\begin{pycode}[Key presses with onkey]
def add_star():
    x = random.randint(-380, 380)
    y = random.randint(20, 280)
    t.penup()
    t.goto(x, y)
    t.dot(8, YELLOW)

screen.onkey(add_star, "space")   # no () after add_star
screen.listen()                   # start listening for keys
\end{pycode}
Click on the turtle window first, then press the key. Ideas: \texttt{"c"} changes the sky colour,
\texttt{"Up"} draws a firework, \texttt{"r"} adds rain.
\end{stretchchallenge}

\begin{reflection}[After the showcase]
The part of my project I am proudest of:
\answerlines{2}
One bug I found and how I fixed it:
\answerlines{2}
\end{reflection}

\end{document}
